\subsection{Свободные колебания полубесконечной струны}

\stepcounter{equation}
\begin{equation}\tag{1}\label{eq:semi-dirichlet}
  \begin{cases}
    u_{tt} - a^2 u_{xx} = 0, & x > 0,\ t > 0, \\[2pt]
    u|_{t=0} = \varphi(x), \quad u_t|_{t=0} = \psi(x), \\[2pt]
    u|_{x=0} = 0 .
  \end{cases}
\end{equation}

Теперь струна занимает полупрямую, а её левый конец закреплён. Нумерация задач в
этом параграфе начинается заново.

\subsubsection*{Метод нечётного продолжения}

Продолжим данные на всю ось нечётно:
\[
  \Phi(x) =
  \begin{cases}
    \varphi(x), & x \ge 0, \\
    -\varphi(-x), & x \le 0,
  \end{cases}
  \qquad\qquad
  \Psi(x) =
  \begin{cases}
    \psi(x), & x \ge 0, \\
    -\psi(-x), & x \le 0 .
  \end{cases}
\]

\inputfig{figures/fig-odd-extension}

Решим на всей прямой задачу с данными $\Phi, \Psi$ по формуле Даламбера:
\[
  u(x,t) = \frac12\Big[\Phi(x-at) + \Phi(x+at)\Big]
         + \frac{1}{2a}\int\limits_{x-at}^{x+at} \Psi(y)\, dy .
\]

\begin{idea}
Ключ: если $\Phi$ и $\Psi$ нечётны, то $u|_{x=0} = 0$ автоматически. Действительно,
при $x=0$
\[
  \frac12\big[\Phi(-at) + \Phi(at)\big] = 0
  \quad\text{(нечётность $\Phi$)},
\]
\[
  \int\limits_{-at}^{at} \Psi(y)\, dy = 0
  \quad\text{(интеграл нечётной функции по симметричному отрезку)} .
\]
Значит построенная функция удовлетворяет и уравнению, и начальным данным (при
$x>0$ они совпадают с исходными), и краевому условию. Условие Дирихле
$u|_{x=0}=0$ «встроено» в симметрию данных.
\end{idea}

\subsubsection*{Явный вид решения}

Если $x - at \ge 0$, то все аргументы попадают на полупрямую $x\ge 0$ и формула
не отличается от случая бесконечной струны. Если же
\[
  x - at < 0,
\]
то работает нечётность:
\[
  \Phi(x-at) = -\varphi(at-x),
  \qquad
  \int\limits_{x-at}^{at-x} \Psi(y)\, dy = 0 ,
\]
поэтому интеграл от $\Psi$ сворачивается до интеграла по отрезку
$[\,at-x,\ x+at\,]$ от исходной $\psi$. Итог:
\[
  \boxed{\
  u(x,t) =
  \begin{cases}
    \dfrac12\Big[\varphi(x-at) + \varphi(x+at)\Big]
      + \dfrac{1}{2a} \displaystyle\int\limits_{x-at}^{x+at} \psi(y)\, dy,
      & x \ge at, \\[18pt]
    \dfrac12\Big[\varphi(x+at) - \varphi(at-x)\Big]
      + \dfrac{1}{2a} \displaystyle\int\limits_{at-x}^{x+at} \psi(y)\, dy,
      & x < at .
  \end{cases}\ }
\]

\begin{idea}
Физический смысл второй строки: при $x < at$ до точки $x$ уже успела дойти волна,
отражённая от закреплённого конца. Слагаемое $-\varphi(at-x)$ — это как раз
отражённая волна, и знак «минус» означает отражение \emph{с переворотом}: у
закреплённого конца горб возвращается впадиной.

\medskip
Прямая $x = at$ в плоскости $(x,t)$ разделяет две зоны: ниже неё отражение ещё не
пришло, выше — уже пришло.
\end{idea}

\subsubsection*{Упражнения}

\begin{exercise}
Решите задачу (свободный конец вместо закреплённого; $\varphi, \psi$ заданы на
полуоси):
\stepcounter{equation}
\begin{equation}\tag{2}\label{eq:semi-neumann}
  \begin{cases}
    u_{tt} - a^2 u_{xx} = 0, & x > 0,\ t > 0, \\[2pt]
    u|_{t=0} = \varphi(x), \quad u_t|_{t=0} = \psi(x), \\[2pt]
    u_x|_{x=0} = 0 .
  \end{cases}
\end{equation}
\end{exercise}

\begin{idea}
Подсказка: условие $u_x|_{x=0} = 0$ (Нейман) получается тем же приёмом, но с
\textbf{чётным} продолжением данных. Отражение тогда идёт без переворота знака.
\end{idea}

\begin{exercise}
Изобразите решения задач~\eqref{eq:semi-dirichlet} и~\eqref{eq:semi-neumann} со
следующими начальными условиями:
\begin{enumerate}[label=(\alph*), leftmargin=2em]
  \item $\varphi(x)$ — «треугольник» с вершиной в точке $x_0 > 0$, $\psi \equiv 0$;
  \item $\varphi \equiv 0$, $\psi(x) = \delta(x - x_0)$.
\end{enumerate}
\end{exercise}

\begin{exercise}
Решите задачи для уравнения теплопроводности на полупрямой:
\[
  \begin{cases}
    u_t - \alpha u_{xx} = 0, & x > 0,\ t > 0, \\
    u|_{t=0} = \varphi(x), \\
    u|_{x=0} = 0 ;
  \end{cases}
  \qquad\qquad
  \begin{cases}
    u_t - \alpha u_{xx} = 0, & x > 0,\ t > 0, \\
    u|_{t=0} = \varphi(x), \\
    u_x|_{x=0} = 0 .
  \end{cases}
\]
\end{exercise}

\begin{idea}
Тот же метод отражения работает и здесь: нечётное продолжение для условия
Дирихле, чётное — для Неймана, а затем свёртка с $\Phi$ из напоминания в начале
конспекта.
\end{idea}
